 \documentclass[]{article}
\usepackage {amsmath}
\usepackage[a4paper, total={7.3in, 10.3in}]{geometry}
\usepackage{cleveref}

\usepackage[titletoc]{appendix}

\usepackage{amssymb}
%\usepackage{xcolor}
\usepackage[x11names, rgb]{xcolor}
\usepackage[utf8]{inputenc}
\usepackage{mathtools}
\usepackage{xcolor}
\usepackage{braket}
\usepackage{ulem}
\usepackage{cancel}
\usepackage{graphicx}
\usepackage{mathtools}
\usepackage{amsmath}  
\usepackage{amsfonts} 
\usepackage{graphicx}
\usepackage{amssymb} 
\usepackage{amsmath}
\usepackage{mathrsfs}
\usepackage{empheq}
\usepackage{amsthm}
 \usepackage{braket}
 \usepackage{amsmath}
\DeclareMathOperator\arctanh{arctanh}
\usepackage[utf8]{inputenc}
\usepackage[english]{babel}
\usepackage{graphicx}
\newtheorem{theorem}{Theorem}[section]
\newtheorem{corollary}{Corollary}[theorem]
\newtheorem{lemma}[theorem]{Lemma}
\numberwithin{equation}{section}

\def\epp{\epsilon^{\prime}}
\def\ep{\epsilon}
\def\vep{\varepsilon}
\def\la{\langle}
\def\ra{\rangle}
\def\ppg{\pi^+\pi^-\gamma}
\def\vp{{\bf p}}
\def\ko{K^0}
\def\kb{\bar{K^0}}
\def\ka{\kappa}
\def\al{\alpha}
\def\ab{\bar{\alpha}}
\def\be{\begin{equation}}
\def\ee{\end{equation}}
\def\bea{\begin{eqnarray}}
\def\eea{\end{eqnarray}}
\def\wh{\widehat}

\usepackage{enumitem}
\title{Response to Referee's Comment}
%\author{\textbf{Hariprashad Ravikumar}} 

\date{}
\begin{document}
	\maketitle

\section*{Referee's Comment}

The authors ......   consistent. 
 The authors indicate a possible extension to the (3+1)-dimensional conformal group with applications to QCD and hadron physics. To strengthen the physical motivation, however, the manuscript would benefit from a more concrete discussion of the specific physical problems this framework is expected to address. 

\noindent
I recommend publication in PRD, subject to clarification on this point. 

\section*{Reply}

We thank the referee for careful reading of our manuscript and providing useful comment to strengthen the physical motivation. We have now added the following paragraph at the end of Summary and Conclusion for a more concrete discussion of the physical problems this framework is expected to address.

\vskip 0.5cm
\noindent
===================\\
As we discussed the correspondence 
 between $(0+1)$d quantum mechanics and AdS$_2$[19] and extended it effectively to the correspondence between $(1+1)$d quantum field theory and AdS$_3$, we look forward to exploring the analysis of the holographic dual correspondence between $(3+1)$d QCD and AdS$_5$, implementing further the interpolation between IFD and LFD. With this analysis, we expect to discuss the effectiveness of the light-front QCD as it maximizes the number of the kinematic generators and provides the link between the first principle QCD and the light-front quark model for the hadron phenomenology. The effectiveness of handling the relativistic interactions among the constituents of the hadrons is highly relevant to the discussion on the advantages that the different forms of relativistic dynamics can offer[36]. 
\vskip 0.5cm

\begin{itemize}
  \item[] [19] V. de Alfaro, S. Fubini, and G. Furlan, Conformal invariance in quantum mechanics, Il Nuovo Cimento A (1965-1970) 34, 569 (1976).
  \item[] [36] C.-R. Ji, Relativistic quantum invariance of QED and QCD, The European Physical Journal Special Topics 10.1140/epjs/s11734-025-01803-9 (2025).
\end{itemize}
===================\\

\noindent
We thank again the referee's useful comment which allow us to improve our manuscript. 

\end{document}